%========================= INTRODUÇÃO =========================
\section{Introdução}

\subsection{Contextualização}

No Brasil, estima-se que mais de 10 milhões de pessoas possuam algum grau de
deficiência auditiva (IBGE, 2023). A Língua Brasileira de Sinais (Libras) é
reconhecida oficialmente como meio legal de comunicação e expressão pela Lei
Federal n.\ 10.436/2002 e regulamentada pelo Decreto n.\ 5.626/2005, que
determina sua inclusão em serviços públicos, na educação e na saúde (BRASIL,
2002; BRASIL, 2005). Na prática, contudo, essa determinação legal é raramente
cumprida de forma ampla, em razão da escassez de intérpretes qualificados frente
à demanda real da população surda.

Essa lacuna se traduz em uma dependência quase exclusiva da presença física de um
intérprete humano para que a comunicação entre pessoas surdas e ouvintes se
estabeleça. Hospitais, bancos, repartições públicas e ambientes educacionais
frequentemente não dispõem desse recurso de forma contínua, de modo que qualquer
atendimento emergencial ou cotidiano se transforma em uma barreira concreta para
o indivíduo surdo.

O avanço recente de técnicas de visão computacional e aprendizado profundo
(\textit{deep learning}) abre uma via alternativa de mediação dessa comunicação:
sistemas capazes de reconhecer automaticamente gestos de língua de sinais a
partir de uma câmera comum, sem necessidade de \textit{hardware} especializado. É
nesse contexto que se insere o presente trabalho.

\subsection{Problema e justificativa}

O problema central que motiva este trabalho pode ser formulado da seguinte
maneira: como reduzir a dependência de um intérprete humano presencial para
viabilizar a comunicação básica entre pessoas surdas e ouvintes, utilizando
apenas equipamentos amplamente acessíveis, como um computador com webcam?

Do ponto de vista social, a justificativa é direta: qualquer solução tecnológica
capaz de mitigar essa barreira, ainda que parcialmente, tem potencial de impacto
imediato sobre o cotidiano da comunidade surda brasileira. Do ponto de vista
acadêmico, a justificativa se apoia em uma lacuna identificada na literatura: a
maior parte dos sistemas de reconhecimento de língua de sinais (\textit{Sign
Language Recognition} — SLR) descritos internacionalmente concentra-se na Língua
Americana de Sinais (ASL) ou em línguas de sinais europeias. Sistemas voltados
especificamente para a Libras, com coleta de dados realizada no Brasil, ainda são
relativamente escassos e representam uma contribuição genuína ao campo, com
potencial direto de impacto social sobre a população brasileira (REZENDE;
ALMEIDA; GUIMARÃES, 2021).

\subsection{Objetivos}

\subsubsection{Objetivo geral}

Desenvolver um sistema funcional de reconhecimento de gestos dinâmicos de Libras
em tempo real, capaz de converter sequências de gestos em texto e formar frases
completas por meio de técnicas de aprendizado profundo e visão computacional.

\subsubsection{Objetivos específicos}

\begin{itemize}[leftmargin=2cm]
  \item implementar um \textit{pipeline} de captura e processamento de vídeo com
  detecção de \textit{landmarks} de mãos via MediaPipe (21 pontos por mão, 30
  FPS);
  \item construir e treinar um modelo híbrido CNN + LSTM capaz de classificar
  sequências temporais de gestos com acurácia superior a 90\%;
  \item criar um sistema de coleta de dados interativo para construção de um
  \textit{dataset} próprio de Libras, com suporte a múltiplos usuários;
  \item implementar a formação automática de frases com \textit{debounce}
  temporal e suporte a três modos de operação: palavras inteiras, soletração
  letra a letra e modo misto;
  \item prover uma interface visual em tempo real que exiba o gesto detectado, o
  percentual de confiança da predição, o progresso do \textit{buffer} de captura
  e a frase formada;
  \item validar o sistema com integrantes da comunidade surda e usuária de
  Libras, registrando métricas de uso e percepção dos participantes.
\end{itemize}

\subsection{Estrutura do trabalho}

Este trabalho está organizado em sete seções, incluindo esta introdução. A
Seção~2 apresenta o referencial teórico, abordando os aspectos linguísticos da
Libras e os fundamentos técnicos de visão computacional e aprendizado profundo
utilizados. A Seção~3 discute os trabalhos relacionados e o estado da arte em
reconhecimento de língua de sinais. A Seção~4 descreve a metodologia adotada,
incluindo o protocolo de validação com usuários. A Seção~5 detalha o sistema
desenvolvido. A Seção~6 apresenta e discute os resultados obtidos, suas
limitações e os trabalhos futuros. Por fim, a Seção~7 traz as considerações
finais.
